% !PRINTED 427
\weberpage{427}
To obtain the \gesperrt{Galois} group of this equation of degree $6$, we have to examine the influence which the application of the icosahedral substitutions has on the system of subsidiary groups (2). For this it suffices to consider the substitutions $\Theta$, $\psi$, $\chi$. But, if we attend to the order of the subsidiary groups in (2), we have by \S.\,58, (23), (25):
\[
\begin{aligned}
P\Theta &= Q + Q\chi\Theta + Q\chi\Theta^2 + Q\chi\Theta^3 + Q\chi\Theta^4 + Q\chi \\
P\psi &= Q + Q\chi + Q\chi\Theta^4 + Q\chi\Theta^3 + Q\chi\Theta^2 + Q\chi\Theta \\
P\chi &= Q\chi + Q + Q\chi\Theta + Q\chi\Theta^3 + Q\chi\Theta^2 + Q\chi\Theta^4,
\end{aligned}
\]
and from this it follows that the indices of $U$ are permuted as follows:
\[
\begin{array}{c@{\quad}cccccc}
 & \infty, & 0, & 1, & 2, & 3, & 4 \\
\cline{2-7}
\Theta) & \infty & 1 & 2 & 3 & 4 & 0 \\
\psi) & \infty & 0 & 4 & 3 & 2 & 1 \\
\chi) & 0 & \infty & 1 & 3 & 2 & 4.
\end{array}
\]
If we denote the index generally by $\xi$, and take $\xi$, as in the seventh chapter, with respect to the modulus 5, so that two numbers congruent with respect to the modulus 5 count as not distinct, then we can represent these permutations thus:
\[
\Theta = (\xi,\ \xi+1),\quad \psi = (\xi,\ -\xi),\quad \chi = \left(\xi,\ \frac{1}{\xi}\right)\!.
\]
\gesperrt{But from this it follows that the group of our equation of degree} $6$ \gesperrt{coincides with the congruence group} $L_5$ (\S.\,68)\gesperrt{.}

To arrive at the formation of resolvents of degree $6$, we must proceed in a way similar to that used for the resolvents of degree $5$.

We have already become acquainted in \S.\,55 with the fundamental forms of the dihedral groups. We now represent them in the form:
\[
\varphi_1 = x_1x_2,\quad \varphi_2 = x_1^5+ix_2^5,\quad \varphi_3 = x_1^5-ix_2^5.
\]
If we take $\Theta$ and $\psi$ with the determinant 1, that is
\[
\Theta = \begin{pmatrix} e^{\frac{\pi i}{5}}, & 0 \\ 0, & e^{-\frac{\pi i}{5}} \end{pmatrix},\quad \psi = \begin{pmatrix} 0, & 1 \\ -1, & 0 \end{pmatrix};
\]
then the fundamental forms $\varphi_1, \varphi_2, \varphi_3$ change in the following way:
\[
\begin{array}{c@{\quad}ccc}
 & \varphi_1 & \varphi_2 & \varphi_3 \\
\cline{2-4}
\Theta) & \varphi_1 & -\varphi_2 & -\varphi_3 \\
\psi) & -\varphi_1 & -i\varphi_2 & i\varphi_3.
\end{array}
\]
